\section*{第 \ref{ch:b1:rifts-ridges-margins} 章——裂谷、洋脊与被动大陆边缘}

\begin{solution}{exo:b1:rifts-ridges-margins:1}
（a）岩墙在平面上覆盖 $\qty{60000}{m}\times\qty{8}{m} = \qty{4.8e5}{m^2}$，
故其平均高度为 $\qty{2.5e9}{m^3}/\qty{4.8e5}{m^2} =
\qty{5200}{m}$，约 \qty{5}{km}：岩墙切穿了\omterm{def:b1:layered-earth:layers}{地壳}的大部分。
（b）$\qty{8}{m}/\qty{0.018}{m.\text{年}^{-1}} = 444$ 年，约四个世纪。
\end{solution}

\begin{solution}{exo:b1:rifts-ridges-margins:2}
（a）\omterm{def:b1:cross-sections:throw}{水平断距}等于\omterm{def:b1:cross-sections:throw}{垂直断距}除以\omterm{def:b1:geological-maps:attitude}{倾角}的正切：
$3.0/\tan\ang{60} = \qty{1.73}{km}$ 和 $0.8/\tan\ang{70} = \qty{0.29}{km}$；
剖面的伸展量为 \qty{2.02}{km}。（b）断层活动前，剖面
原宽 $25 - 2.02 = \qty{22.98}{km}$，约 \qty{23}{km}：此后已伸长约
\qty{9}{\percent}。
\end{solution}

\begin{solution}{exo:b1:rifts-ridges-margins:3}
超慢速扩张（低于 \qty{20}{mm.\text{年}^{-1}}）：加克尔洋脊（12）、红海中部（15）
以及雷克雅内斯洋脊（19），后者接近界限。慢速扩张（20 至
\qty{55}{mm.\text{年}^{-1}}）：卡尔斯伯格洋脊西段（23）和南大西洋
（33）。快速扩张（高于 \qty{75}{mm.\text{年}^{-1}}）：东太平洋海隆（143）。
超慢速与慢速洋脊应在正断层之间发育轴向谷；
东太平洋海隆则应为轴部隆起。
\end{solution}

\begin{solution}{exo:b1:rifts-ridges-margins:4}
（a）$\qty{3.0}{km^2}\times\qty{7.1}{km} = \qty{21.3}{km^3}$，约
每年 \qty{21}{km^3}。（b）$\qty{3.028e8}{km^2}/\qty{3.0}{km^2.\text{年}^{-1}} =
\qty{1.01e8}{\text{年}}$，即 \qty{101}{Ma}。最老的洋底，约 \qty{190}{Ma}，
之所以更老，是因为洋底并非按年龄顺序被消亡：俯冲会吞没
各个年龄的洋底，部分老洋底因未被海沟波及而得以保留，
且生成速率曾发生变化。
\end{solution}

\begin{solution}{exo:b1:rifts-ridges-margins:5}
（a）$2500 + 350\times2 = \qty{3200}{m}$，$2500 + 350\times6 = \qty{4600}{m}$，
$2500 + 350\times8 = \qty{5300}{m}$。（b）$t = \bigl((d - 2500)/350\bigr)^2$：
$(2100/350)^2 = \qty{36}{Ma}$ 和 $(2800/350)^2 = \qty{64}{Ma}$。
（c）洋底在 4 至 \qty{36}{Ma} 之间下沉 \qty{1400}{m}，而
在 \qty{36}{Ma} 至 \qty{64}{Ma} 之间仅下沉 \qty{700}{m}，历时相近。
\end{solution}

\begin{solution}{exo:b1:rifts-ridges-margins:6}
有效热膨胀系数为
$\num{3.2e-5}\times3300/2300$ $= \qty{4.591e-5}{K^{-1}}$，且
$\sqrt{\kappa\times\qty{1}{Ma}}$ $= \sqrt{\num{e-6}\times\num{3.156e13}}$ $=
\qty{5618}{m}$。于是 $C = (2/\sqrt\pi)\times\num{4.591e-5}\times1350\times5618$ $=
1.1284\times\num{4.591e-5}\times1350\times5618$ $= \qty{393}{m.Ma^{-1/2}}$。对于
温度低 \qty{250}{K} 的\omterm{def:b1:layered-earth:layers}{地幔}，$C$ 与 $T_m - T_0$ 成正比：
$393\times1100/1350 = \qty{320}{m.Ma^{-1/2}}$。较冷洋盆的洋底
下沉更慢，遵循较小的常数。本章拟合得到的 \num{350}
并非此类情形：它落在参数的不确定度范围之内。
\end{solution}

\begin{solution}{exo:b1:rifts-ridges-margins:7}
（a）$\qty{e-6}{cal.cm^{-2}.s^{-1}} = \num{4.184e-6}\times\num{e4} =
\qty{4.184e-2}{W.m^{-2}} = \qty{41.84}{mW.m^{-2}}$，故常数为
$11.3\times41.84 = \qty{473}{mW.m^{-2}.Ma^{1/2}}$，比
\cref{ch:b1:heat-budget}中的 490 低 \qty{3.5}{\percent}，两条拟合结果一致。（b）$q = 473$，
$473/\sqrt{10} = 150$ 和 $473/\sqrt{150} = \qty{38.6}{mW.m^{-2}}$。
平方根律持续下降，而板块模型趋于一个常值：
在 \qty{150}{Ma} 处平方根律的数值偏低，因为\omterm{thm:b1:rifts-ridges-margins:depth-age}{深度–年龄关系}
下降过深。
\end{solution}

\begin{solution}{exo:b1:rifts-ridges-margins:8}
设北段位于 $x = 0$，南段位于 $x = -40$
（单位 km），\omterm{def:b1:rifts-ridges-margins:transform}{转换断层}位于两者之间，即 $y = 0$。（a）紧邻\omterm{def:b1:rifts-ridges-margins:transform}{转换断层}以北，洋底
位于北段以西并向西运动；紧邻其以南处，
则位于南段以东并向东运动。从任一侧看，
远侧向左运动：此为左行错动。（b）
洋脊看上去被右行错开，其南段位于
以西：方向相反，正如\cref{prop:b1:rifts-ridges-margins:transform}
所预测。（c）地震只发生在两洋脊端点之间的 \qty{40}{km} 范围内，
那里两侧分属不同板块；在此之外，沿
\omterm{def:b1:rifts-ridges-margins:transform}{破碎带}，两侧同属一个板块并共同运动。
\end{solution}

\begin{solution}{exo:b1:rifts-ridges-margins:9}
半速率为 \qty{10.6}{mm.\text{年}^{-1}}，即 \qty{10.6}{km.Ma^{-1}}。（a）年龄
阶跃为 $350/10.6 = \qty{33.0}{Ma}$。（b）在\omterm{def:b1:rifts-ridges-margins:transform}{转换断层}末端，
\qty{33.0}{Ma} 的洋底（$2500 + 350\sqrt{33.0} = \qty{4511}{m}$）与
年龄为 0 的洋底（\qty{2500}{m}）相对，阶跃约为 \qty{2010}{m}。再向外 \qty{100}{km}，
两侧年龄分别为 $450/10.6 = 42.5$ 与 $100/10.6 = \qty{9.4}{Ma}$，深度分别为
4781 与 \qty{3575}{m}，阶跃约 \qty{1210}{m}。再向外 \qty{1000}{km}，
年龄分别为 127.4 与 \qty{94.3}{Ma}，深度分别为 6450 与 \qty{5899}{m}：
\qty{550}{m}。（这些年龄超过 \qty{80}{Ma}，平方根律在此已
下降过深；板块律给出的阶跃会更小。）
\end{solution}

\begin{solution}{exo:b1:rifts-ridges-margins:10}
（a）各年龄所占面积相等时，平均深度为 $d(t)$ 在
$[0, A]$ 上的平均值：
\[
  \frac1A\int_0^A\bigl(d_r + C\sqrt t\bigr)\mathrm dt = d_r + \frac CA\cdot\frac23A^{3/2}
  = d_r + \frac23C\sqrt A.
\]
（b）当 $A = \qty{180}{Ma}$ 时：$2500 + \frac23\times350\times13.416 =
\qty{5630}{m}$；当 \qty{120}{Ma} 时：$2500 + \frac23\times350\times10.954 =
\qty{5056}{m}$，浅 \qty{574}{m}。（c）设海平面上升 $r$，
负载洋底下沉 $s = r\rho_w/(\rho_m - \rho_w)$。水的体积保持不变，
在同一面积上须填充相同深度，故 $r + s = \qty{574}{m}$，
$r\,\rho_m/(\rho_m - \rho_w) = 574$，且 $r = 574\times2300/3300 =
\qty{400}{m}$，仅由扩张速率一项便可产生数百米的变化。
\end{solution}

\begin{solution}{exo:b1:rifts-ridges-margins:11}
（a）$\sqrt{x - a} = \sqrt x\sqrt{1 - a/x} \approx \sqrt x\,(1 - a/2x)$，当
$a \ll x$ 时，故
$\Delta d \approx C\sqrt{x/u}\cdot a/(2x) = Ca/(2\sqrt{ux})$.
（b）当 $x = \qty{450}{km}$ 时：精确值为 \qty{1205}{m}，近似值
为 $350\times350/(2\sqrt{10.6\times450}) = \qty{887}{m}$；当 \qty{1350}{km} 时：
\qty{550}{m} 与 \qty{512}{m}；当 \qty{3350}{km} 时：\qty{334}{m} 与
\qty{325}{m}。近似阶跃与精确阶跃之比仅取决于
$r = a/x$，即 $(r/2)/(1 - \sqrt{1 - r})$，当 $r = 0.36$ 时其值为 0.9。
该近似的误差在 \qty{10}{\percent} 以内，条件是超过 $x = 350/0.36 \approx
\qty{970}{km}$，约为离洋脊三个错距。
\end{solution}

\begin{solution}{exo:b1:rifts-ridges-margins:12}
（a）$d_2 - d_1 = 3671.1 - 3606.8 = \qty{64.3}{m}$，当 \qty{10}{Ma} 时；
$4072.8 - 4065.3 = \qty{7.5}{m}$，当 20 时；$4415.4 - 4417.0 = \qty{-1.7}{m}$，当
30 时；$4707.5 - 4713.6 = \qty{-6.1}{m}$，当 40 时。（b）$d_2 - d_1$ 在 20 时为正，
在 30 时为负，两定律均连续，故二者在其间必有相等之处。
二分法：当 25 时，$d_2 - d_1 = +0.9$；当 27.5 时，$-0.7$；当 26.25 时，$+0.01$；当
26.875 时，$-0.3$；当 26.56 时，$-0.2$；当 26.41 时，$-0.08$；当 26.33 时，$-0.04$：
根位于 26.25 与 26.33 之间，在 \qty{26.3}{Ma} 处，二者均给出
\qty{4294}{m}。（c）$d_1' = 175/\sqrt{26.3} = \qty{34.1}{m.Ma^{-1}}$ 与
$d_2' = (3200/62.8)\,e^{-26.3/62.8} = \qty{33.5}{m.Ma^{-1}}$。两定律在该处的取值与斜率几乎相同，
故拼接后的曲线在此处没有
可见的拐折。
\end{solution}

\begin{solution}{pb:b1:rifts-ridges-margins:1}
\textbf{1.} \omterm{def:b1:cross-sections:throw}{水平断距} $= $ \omterm{def:b1:cross-sections:throw}{垂直断距}$/\tan$（\omterm{def:b1:geological-maps:attitude}{倾角}）：F1 \qty{346}{m}，F2 \qty{770}{m}，
F3 \qty{2014}{m}，F4 \qty{163}{m}，F5 \qty{462}{m}，F6 \qty{1172}{m}。F3 的\omterm{def:b1:cross-sections:throw}{垂直断距}和\omterm{def:b1:cross-sections:throw}{水平断距}最大，
为边界断层。

\textbf{2.} 伸展量为各断层\omterm{def:b1:cross-sections:throw}{水平断距}之和，即 \qty{4927}{m}，约 \qty{4.9}{km}。裂谷
在断裂前宽 $60 - 4.9 = \qty{55.1}{km}$。

\textbf{3.} \omterm{def:b1:geological-maps:attitude}{走向}为 \ang{020} 的裂谷，其法线指向 \ang{110}；
运动朝向 \ang{105}，与之成 \ang{5} 夹角。其垂直分量
为 $6.0\cos\ang{5} = \qty{5.98}{mm.\text{年}^{-1}}$。

\textbf{4.} $\qty{4.927}{km}/\qty{5.98}{km.Ma^{-1}} = \qty{0.82}{Ma}$.

\textbf{5.} 在 \qty{3.2}{Ma} 内，两板块相互分离的距离为
$5.98\times3.2 = \qty{19.1}{km}$，跨越裂谷。断层
承担其中的 $4.93/19.1 = \qty{26}{\percent}$。

\textbf{6.} 其余约四分之三，即 \qty{14.2}{km}，不在
本节所述断层之中。它们或许由沿轴线侵入的岩墙承担，此类岩墙
可在不使裂谷断裂的情况下使其张开，如阿法尔所示；或由本节之外的断层承担。
若岩墙宽 \qty{8}{m}，则 $\qty{14200}{m}/\qty{8}{m} \approx 1780$ 条岩墙
需要，平均每隔 $3.2\times10^6/1780 \approx 1800$ 年一条。

\textbf{7.} 由 C2An 得：$(21.0 + 20.2)/2.581 = 41.2/2.581 =
\qty{15.96}{mm.\text{年}^{-1}}$。由 C3n 得：$(41.0 + 40.2)/5.235 = 81.2/5.235 =
\qty{15.51}{mm.\text{年}^{-1}}$.

\textbf{8.} $(41.2 - 2.2)/2.581 = \qty{15.11}{mm.\text{年}^{-1}}$，与
$(81.2 - 2.2)/5.235 = \qty{15.09}{mm.\text{年}^{-1}}$：修正后二者在
\qty{0.02}{mm.\text{年}^{-1}} 之内相符。修正对较年轻的磁\omterm{def:b1:plate-kinematics:magnetic-anomaly}{极性时}影响更大，其
距离较短。

\textbf{9.} 向外的位移使洋脊一翼的两个边缘
移动相同的量，因此在二者之差中相互抵消。对两翼合并计算，
$(81.2 - 41.2)/(5.235 - 2.581) = 40.0/2.654 = \qty{15.07}{mm.\text{年}^{-1}}$.

\textbf{10.} 东北翼：$(41.0 - 21.0)/2.654 = \qty{7.54}{mm.\text{年}^{-1}}$；
西南翼：$(40.2 - 20.2)/2.654 = \qty{7.54}{mm.\text{年}^{-1}}$。扩张是
对称的；两翼之间 \qty{0.8}{km} 的差值对两个磁\omterm{def:b1:plate-kinematics:magnetic-anomaly}{极性时}相同，
是一个恒定偏移，例如轴线被放置在偏离
中心 \qty{0.4}{km} 处。

\textbf{11.} $\cos\Delta = \sin\ang{30.9}\sin\ang{18.0} +
\cos\ang{30.9}\cos\ang{18.0}\cos\ang{16.4} = 0.15870 + 0.78287 = 0.94157$，故
$\Delta = \ang{19.69}$。取 $\omega R = 0.403\times(\pi/180)\times6371 =
\qty{44.81}{mm.\text{年}^{-1}}$，则速率为 $44.81\sin\ang{19.69} =
\qty{15.10}{mm.\text{年}^{-1}}$，即第 8、9 题的修正速率。

\textbf{12.} $\cos\Delta = \sin\ang{30.9}\sin\ang{26.0} +
\cos\ang{30.9}\cos\ang{26.0}\cos\ang{11.4} = 0.22512 + 0.75600 = 0.98112$，故
$\Delta = \ang{11.15}$，速率为 $44.81\sin\ang{11.15} =
\qty{8.7}{mm.\text{年}^{-1}}$.

\textbf{13.} 预测值 15.1 与 \qty{8.7}{mm.\text{年}^{-1}} 与
测得的 15 与 \qty{9}{mm.\text{年}^{-1}} 相符。速率为 $\omega R\sin\Delta$，
\omterm{def:b1:plate-kinematics:euler-pole}{欧拉极}位于红海西北方，在利比亚东北部：越靠南，离\omterm{def:b1:plate-kinematics:euler-pole}{欧拉极}越远，
张开就越快。

\textbf{14.} $15.09\times5.235 = \qty{79}{km}$ 的洋底。

\textbf{15.} $5000/120.6 = \qty{41}{km.Ma^{-1}}$，即 \qty{41}{mm.\text{年}^{-1}}，平均
而言。

\textbf{16.} 按中部速率，$5000/15.1 = \qty{331}{Ma}$；
$5000/8.7 = \qty{575}{Ma}$，按北部速率。

\textbf{17.} 绕附近的\omterm{def:b1:plate-kinematics:euler-pole}{欧拉极}张开，红海如一把
剪刀般张开：其宽度与 $\sin\Delta$ 成正比，始终是
朝\omterm{def:b1:plate-kinematics:euler-pole}{欧拉极}方向收窄的楔形。以 \qty{8.7}{mm.\text{年}^{-1}} 计，北部需要
$79/8.7 = \qty{9.1}{Ma}$ 方能形成中部
现有的 \qty{79}{km} 洋底。

\textbf{18.} 以 \qty{15}{mm.\text{年}^{-1}} 计，红海轴线属\omterm{def:b1:rifts-ridges-margins:ridge-classes}{超慢速扩张洋脊}；
预计应有由正断层夹峙的轴向谷。

\textbf{19.} $2500 + 350\times1 = \qty{2850}{m}$;
$2500 + 350\sqrt{10} = \qty{3607}{m}$; $2500 + 350\sqrt{40} = \qty{4714}{m}$.

\textbf{20.} $t = (1500/350)^2 = \qty{18.4}{Ma}$.

\textbf{21.} 在最初 \qty{10}{Ma} 内，其下沉 $350\sqrt{10} = \qty{1107}{m}$；
在 70 至 \qty{80}{Ma} 之间，$350(\sqrt{80} - \sqrt{70}) = 350\times(8.944 -
8.367) = \qty{202}{m}$，同样时长内约为前者的五分之一。

\textbf{22.} $\mathrm dd/\mathrm dt = C/(2\sqrt t) = 175/\sqrt{80} =
\qty{19.6}{m.Ma^{-1}}$，约合每年 \qty{0.02}{mm}。

\textbf{23.} $490/\sqrt{80} = \qty{54.8}{mW.m^{-2}}$.

\textbf{24.} 平方根定律：$2500 + 350\sqrt{80} = 2500 + 350\times8.944 =
\qty{5630}{m}$。板块定律：$6400 - 3200\,e^{-80/62.8} = 6400 - 3200\times0.2798
= \qty{5505}{m}$。两者相差 \qty{126}{m}。在 \qty{80}{Ma} 时两定律仍均适用：
平方根定律适用至约 \qty{80}{Ma}，板块定律自约
\qty{20}{Ma} 起适用；此后平方根定律给出的深度偏深。

\textbf{25.} 平方根定律给出 \qty{5630}{m}：精确到十分之
一千米，新生洋底在年龄为
8000 万年时的深度为 \textbf{\qty[text-series-to-math=true]{5.6}{km}}
（板块定律给出 \qty{5.5}{km}）。
\end{solution}
